\clearpage
\fancyhead[L]{\small ZACH. L'ATTRACTION DES MONTAGNES.}
\hypertarget{native-574-lower}{}
\begin{center}\small Göttingische gelehrte Anzeigen. Stück 210. Seiten 2089..2100. 31. Dezember 1814.\\[3pt]\rule{.72\linewidth}{.3pt}\end{center}

\emph{L'attraction des montagnes, et ses effets sur les fils à plomb ou sur les niveaux des
instrumens d'astronomie, constatés et determinés par des observations astronomiques et
géodésiques faites en} 1810, \emph{à l'ermitage de notre dame des anges, sur le mont de Mimet, et
au fanal de l'isle de Planier près de Marseille; suivis de la description géométrique de la
ville de Marseille et de son territoire. Par le baron} \textsc{de Zach}. 1814.
Avignon. Bei \textsc{Seguin} dem Älteren. XX und 716 Seiten im Großoktavformat.

Die Beobachtungen, die \textsc{Bouguer} und \textsc{Condamine} im Jahr 1738 am Chimborasso und \textsc{Maskelyne} im Jahr 1774 am Shehallien in Schottland zum Einfluss großer Gebirgsmassen auf die Richtung der Schwerkraft angestellt hatten, waren bisher die einzigen direkten Versuche zu diesem interessanten Gegenstand. Die ersteren hatten für die Ablenkung des Lots von seiner natürlichen Richtung $7''5$ ergeben, die anderen $5''8$. Streng genommen sind aber nur die Beobachtungen in Schottland eindeutig zuverlässig. Die Beobachtungen in Peru sind es keineswegs, da sie unter zu ungünstigen Umständen und mit verhältnismäßig zu unvollkommenen Instrumenten ausgeführt wurden. Angesichts der Bedeutung des Gegenstands und der inzwischen so großen Fortschritte der praktischen Astronomie ist es daher höchst wünschenswert, dass weitere solche interessanten Versuche angestellt werden. Die vorliegende Arbeit eines der führenden Beobachter ist somit als echter Gewinn für die Wissenschaft anzusehen. Das hier besprochene Werk enthält sämtliche zu diesem Unternehmen gehörigen Beobachtungen bis ins kleinste Detail, stellt die angewandten Beobachtungs- und Rechenverfahren ausführlich dar und bietet außerdem zahlreiche sorgfältig berechnete Hilfstafeln, die sich auch anderswo mit großem Vorteil nutzen lassen. Darüber hinaus enthält es eine Fülle verschiedenartiger praktischer Bemerkungen und Urteile sowie gelehrte Notizen und Nebenuntersuchungen. Diesen Vorzug ist man aus den Schriften des Verfassers bereits gewohnt.

\SourceLeaf{575}In der Einleitung führt der Verfasser den Leser zunächst zu dem Standpunkt, von dem aus man das Phänomen der Anziehung der Berge und die Beobachtungen betrachten muss, durch die es sichtbar wird. Dann berichtet er von den einschlägigen bisherigen Erfahrungen und gibt schließlich einen allgemeinen Überblick über seine eigenen Arbeiten. Hier mag es genügen, daran zu erinnern, dass die Schwerkraft nichts anderes ist als die Gesamtwirkung, die alle Bestandteile des Erdkörpers nach dem allgemeinen Anziehungsgesetz auf die Körper an der Erdoberfläche ausüben, verändert durch die Zentrifugalkraft infolge der Erdrotation. Die astronomische Polhöhe, also der Winkel zwischen der Richtung der Schwerkraft und der Äquatorebene, muss daher vom Äquator zum Pol hin regelmäßig und nach Gesetzen zunehmen, die man aus dem Vergleich der verschiedenen Gradmessungen ableitet, sofern man die Erde als regelmäßigen Körper betrachtet. Diese regelmäßige Zunahme wird jedoch dort gestört, wo große Abweichungen der Erdoberfläche von der Normalgestalt merkliche Ablenkungen der Schwerkraft von ihrer natürlichen Richtung hervorrufen. Betrachten wir zwei Orte. An einem von ihnen beeinflusst eine nahe Bergmasse die Richtung der Schwerkraft von Norden nach Süden oder von Süden nach Norden. Am anderen Ort ist der Einfluss derselben Masse entweder nicht mehr merklich oder wirkt in entgegengesetzter Richtung. Die astronomische Beobachtung ergibt dann für den Unterschied der Polhöhen beider Orte einen anderen Wert als die Rechnung aus ihrer gegenseitigen Lage und Entfernung. An dieser Abweichung erkennen wir im ersten Fall die Größe der Ablenkung, im anderen die Summe beider Ablenkungen. Offensichtlich ist also der zweite Fall, der bei \textsc{Maskelyne}s Beobachtungen vorlag, der günstigere. Im ersten Fall ist die Aufgabe doppelt so schwierig, da die ganze Ablenkung immer nur sehr wenige Bogensekunden beträgt und die Beobachtungen daher von größter Genauigkeit sein müssen. Das Gelände, in dem Herr \textsc{von Zach} seine Arbeit ausführte, ließ die Anwendung jenes günstigeren Verfahrens wahrscheinlich nicht zu. Einer seiner beiden Beobachtungspunkte lag am Südhang des Mimet, des höchsten Berges (400 Toisen über der Meeresoberfläche) in einer von Osten nach Westen verlaufenden Reihe von Kalkbergen etwa zwei Meilen nördlich von Marseille. Der andere war der Leuchtturm auf der sehr kleinen Insel Planier, ein paar Meilen südwestlich von Marseille. Am ersten Ort, in 250 Toisen Höhe beim verfallenen Kloster Notre-Dame des Anges, musste der Berg die Richtung der Schwerkraft nach Norden ablenken und dadurch die Polhöhe verringern. Am anderen Ort konnte der Einfluss der Berganziehung wegen der großen Entfernung als unmerklich gelten. Die beobachteten astronomischen Polhöhen mussten also eine kleinere Differenz ergeben als die Rechnung aus der durch geodätische Messungen bestimmten Lage beider Orte, wenn man dabei die regelmäßige Zunahme voraussetzt. Zur Bestimmung der Polhöhen verwendet man gern an beiden Orten dieselben Sterne. Der Unterschied der Polhöhen, auf den es hier allein ankommt, wird dadurch unabhängig von der absoluten Richtigkeit der Sterndeklinationen. Man kann sogar die absoluten Polhöhen außer Betracht lassen und statt ihrer Differenz die Differenz der beobachteten Zenitdistanzen der Sterne verwenden, nachdem man diese von den kleinen periodischen Ungleichheiten befreit und auf denselben Zeitpunkt reduziert hat.

Von den acht Abschnitten, in die Herr \textsc{von Zach} das vorliegende Werk gegliedert hat, enthalten die ersten beiden sämtliche astronomischen Beobachtungen bei Notre-Dame des Anges und auf der Insel Planier. Es waren Beobachtungen dreierlei Art. Man beobachtete die Zenitdistanzen der drei Sterne $\alpha$ Ophiuchus, $\zeta$ Adler und $\alpha$ Adler, ferner Pulversignale zur Bestimmung der Längenunterschiede gegenüber der Sternwarte in Marseille und schließlich Azimute zur Orientierung des Dreiecksnetzes, das die beiden Beobachtungsorte miteinander verband. Die Zeiten \SourceLeaf{576}bestimmte man mit drei Chronometern von \textsc{Emery} mithilfe korrespondierender Sonnenhöhen. Die Zenitdistanzen beobachtete man mit einem zwölfzölligen Vervielfältigungskreis von \textsc{Reichenbach}, die Azimute mit einem achtzölligen Theodoliten desselben Instrumentenbauers. Für Beobachter, die mit Vervielfältigungskreisen arbeiten, ist die hier mitgeteilte allgemeine Reduktionstafel für die außerhalb der Kulmination beobachteten Zenitdistanzen wertvoll. Sie ist genauer und vollständiger berechnet als die Tafeln, die man anderswo findet. Der dritte Abschnitt enthält die terrestrischen Messungen. Eine Basis von 1182,4 Toisen wurde auf der Straße von Marseille nach Aix mit hölzernen Messstangen gemessen, zu deren geringer Veränderlichkeit der Verfasser hier bemerkenswerte Beobachtungen mitteilt. Die horizontalen Winkel der sieben Dreiecke des Netzes, das die beiden Hauptpunkte mit der Basis verband, wurden ebenfalls mit dem Theodoliten beobachtet, meist zehnmal, einige auch öfter. Die größte Abweichung der Summe der drei Winkel in einem Dreieck betrug $4''9$. Im vierten Abschnitt leitet der Verfasser aus diesen geodätischen Messungen den Unterschied der Breite und der Länge zwischen den beiden Beobachtungsorten ab. Dabei legt er die sphäroidische Gestalt der Erde, die Abplattung $\frac{1}{310}$ und den Halbmesser des Erdäquators von 3271604 Toisen zugrunde und verwendet die Formeln, die \textsc{Delambre} in den Méthodes analytiques pour la détermination d'un arc du méridien angegeben hat. Der Breitenunterschied ergibt sich zu $12'\,3''11$, der Längenunterschied zu $15'\,46''46$. Auch für diese Rechnungen teilt der Verfasser verschiedene Hilfstafeln mit. Im fünften Abschnitt werden die Rechenergebnisse nun mit den Ergebnissen der astronomischen Beobachtungen verglichen. Für den Unterschied der Polhöhen ergaben die Beobachtungen die folgenden Werte.

\begin{center}\small\begin{tabular}{@{}l@{\hspace{9pt}}r@{}}
598 Beobachtungen von $\alpha$ Ophiuchus $.\;.\;.\;.$ & $12'\,0''84$ \\
518 Beobachtungen von $\zeta$ Adler $.\;.\;.\;.\;.$ & $12\phantom{'}\,1.26$ \\
654 Beobachtungen von $\alpha$ Adler $.\;.\;.\;.\;.$ & $12\phantom{'}\,1.30$ \\
\end{tabular}\end{center}
\noindent Im Mittel ergaben also 1770 Beobachtungen $12'\,1''13$, sodass $1''98$ als das endgültige Ergebnis für die Wirkung der Anziehung des Mimet am Beobachtungsort Notre-Dame des Anges anzusehen ist. Aus den Beobachtungen der Pulversignale ergaben sich die folgenden Längenunterschiede zwischen der Sternwarte in Marseille und den beiden Orten.

\begin{center}\small\begin{tabular}{@{}l@{\hspace{9pt}}r@{}}
Notre-Dame des Anges $.\;.\;.\;.\;.\;.\;.\;.$ & $7'\,29''25$ \\
Leuchtturm auf der Insel Planier $.\;.\;.\;.$ & $8\phantom{'}\,5.70$ \\
\end{tabular}\end{center}
\noindent Der erste Wert beruht auf 63 Beobachtungen an 11 verschiedenen Tagen, der zweite auf 53 Beobachtungen an 12 Tagen (wobei fünf Beobachtungen eines Tages allein deshalb ausgeschlossen wurden, weil ihr Ergebnis, $7'\,38''1$, von den übrigen zu stark abwich). Der gesamte Längenunterschied zwischen Notre-Dame des Anges und der Insel Planier beträgt nach diesen Beobachtungen also $15'\,34''95$, oder $15'\,35''79$, wenn man ihn auf die Punkte reduziert, auf die sich die geodätischen Messungen beziehen. Er ist somit um $10''67$ kleiner als nach den geodätischen Messungen oder, wenn man gar keine Beobachtung ausschließt, um $13''05$. Man ist jedoch keineswegs berechtigt, auch diesen Unterschied der Anziehung des Berges zuzuschreiben. Vielmehr ist er wenigstens größtenteils auf unvermeidliche Beobachtungsfehler bei der astronomischen Bestimmung des Längenunterschieds zurückzuführen, die von der Zeit abhing und bei weitem nicht den Genauigkeitsgrad der Beobachtungen der Zenitdistanzen erreichen konnte. Im Übrigen war die ganze Längenbestimmung nur eine untergeordnete, für die Hauptaufgabe keineswegs wesentliche Arbeit. Es genügt schon, dass wir daraus lernen, welchen Fehlern bei dem angewandten Verfahren selbst der geübteste Beobachter noch ausgesetzt ist. Ähnliches lehren die Vergleiche der verschiedenen beobachteten Azimute. Überträgt man die auf der Insel Planier bestimmten Azimute auf den Punkt Notre-Dame des Anges, so ergeben sich zwischen \SourceLeaf{577}den so abgeleiteten und den dort unmittelbar beobachteten Azimuten Unterschiede von 7 und 21 Bogensekunden. Diese Abweichungen lassen sich gut erklären, wenn man sie angemessen auf die Azimutbestimmungen und die terrestrischen Winkelbeobachtungen verteilt. Wir kommen jedoch nicht umhin, hier einen kleinen Umstand zu erwähnen, der gerade bei Vergleichen \emph{dieser} Art nicht unwichtig ist. Uns ist aufgefallen, dass beim Zentrieren der Winkel die Abstände der Beobachtungsorte von den Dreieckspunkten, die der Verfasser in Zehntausendsteln einer Toise angibt, ins metrische Maß umgerechnet fast durchweg nur Zehntelmeter, also ganze Dezimeter, ergeben. Das scheint kein Zufall zu sein. Der Verfasser scheint diese Abstände mit einem metrischen Maßstab gemessen und kleinere Teile unbeachtet gelassen zu haben. Bei mehreren Winkeln verändert jedoch ein Fehler von einem Dezimeter (etwa 4 Zoll) die Reduktion um $3''$. Für den Hauptzweck des Verfassers ist das freilich sehr unbedeutend. Will man jedoch durch die Übertragung eines Azimuts auf einen anderen, entfernten Ort und den Vergleich mit einem dort unmittelbar beobachteten Azimut neue Aufschlüsse über die Gestalt der Erde gewinnen, so muss man auch bei dieser Zentrierung umso sorgfältiger vorgehen, je größer die Zahl der Zwischendreiecke ist. Der Verfasser selbst empfiehlt mit Nachdruck, eine solche Arbeit mit Instrumenten von \textsc{Reichenbach} auszuführen, und schlägt dafür die Liparischen Inseln vor. Wir wünschen uns nichts mehr, als dass der Verfasser dieses Unternehmen selbst ausführt. Darauf dürfen wir umso mehr hoffen, als er seinen bisherigen Aufenthaltsort gewechselt hat.

Der sechste Abschnitt ist der gehaltvollste des ganzen Werks. Der Verfasser untersucht zunächst, wie sich mögliche Fehler der einzelnen Arbeiten auf das Endergebnis auswirken können. Es ist leicht einzusehen, dass die Beobachtungen der Zenitdistanzen bei weitem der empfindlichste Teil des Ganzen sind. Angesichts des außergewöhnlich kleinen Endergebnisses, dessen Zuverlässigkeit bei diesen Beobachtungen eine Genauigkeit bis auf Bruchteile von Bogensekunden voraussetzt, könnte mancher umso eher Bedenken haben, als der Verfasser selbst früher bei anderer Gelegenheit erklärt hat, dass alle mit Repetitionskreisen bestimmten absoluten Polhöhen stets mit einer Ungewissheit von mehreren Bogensekunden behaftet bleiben. Diese Behauptung beruhte freilich auf der Erfahrung, dass solche mit \emph{verschiedenen} Kreisen ausgeführten Bestimmungen derartige Unterschiede zeigten, während die mit jedem einzelnen Kreis erzielten Ergebnisse untereinander ausgezeichnet übereinstimmten. Bei der vorliegenden Arbeit wurden dagegen alle Beobachtungen mit ein und demselben Instrument angestellt. Sie stimmen untereinander vorzüglich überein, und hier kommt es allein auf kleine Unterschiede an, unabhängig von den absoluten Bestimmungen selbst. Zur weiteren Bestätigung seiner Ergebnisse nutzt Herr \textsc{von Zach} die zahlreichen astronomischen Beobachtungen, die er teils an der Königlichen Sternwarte in Marseille, teils an seinen eigenen Sternwarten La Capellette und St.\ Peyre angestellt hat. Die Unterschiede zwischen den so bestimmten Polhöhen dieser drei Punkte (die durch die Anziehung der Berge nicht mehr merklich beeinflusst werden) und der Polhöhe der Insel Planier stimmen äußerst genau mit den Ergebnissen der geodätischen Messungen überein, durch die er sie mit seinem Dreiecksnetz verbunden hat. Ein ähnlicher Vergleich mit Notre-Dame des Anges führt dagegen wieder mit recht guter Übereinstimmung zu jenem Unterschied von zwei Bogensekunden. Man könnte vielleicht bezweifeln, ob diese Bestätigung zulässig ist, da sich die fünf Polhöhen auf verschiedene Sterne stützen (Notre-Dame des Anges und die Insel Planier auf $\alpha$ Adler, die drei anderen hingegen auf $\alpha$ und $\beta$ im Kleinen Bären). Wenn man aber gebührend berücksichtigt, dass die verwendeten Deklinationen genau diejenigen sind, die Herr \textsc{von Zach} selbst mit eben diesem Instrument an einem nicht sehr viel weiter nördlich gelegenen Ort, in Mailand, bestimmt hat, sieht man leicht, dass dieser Zweifel unbegründet ist, sofern man nur die Voraussetzung gelten lässt, dass \SourceLeaf{578}etwaige Fehler bei den mit dem Kreis bestimmten absoluten Zenitdistanzen für bestimmte Zenitdistanzen unveränderlich und für wenig verschiedene Zenitdistanzen nahezu gleich sind. Wie man auch über die Ursachen der bemerkenswerten Unterschiede urteilen mag, die Herr \textsc{von Zach} mit \emph{verschiedenen} Kreisen gefunden hat, so ist doch nicht zu leugnen, dass sich die beständige Übereinstimmung der mit demselben Kreis gewonnenen Ergebnisse ohne die genannte Voraussetzung nicht erklären lässt. Aus all diesen Gründen darf man annehmen, dass der gefundene Unterschied von 2 Bogensekunden die Anziehung des Mimet tatsächlich sehr genau wiedergibt. Doch wird kein praktischer Astronom bestreiten, dass bei diesem Ergebnis immer noch eine Ungewissheit von einer halben Bogensekunde (wenn nicht mehr) bleibt und sich das Ergebnis also nicht enger eingrenzen lässt als auf $1''5$ bis $2''5$. Man muss daher bedauern, dass eine so schöne und so sorgfältig ausgeführte Arbeit nur eine so kleine Größe ergeben hat, im Verhältnis zu der die unvermeidliche Ungewissheit so beträchtlich ist. Umso bedauerlicher ist es, dass der zweite Beobachtungsort auf der Insel Planier gewählt wurde und nicht auf der Nordseite unmittelbar am Mimet. Herr \textsc{von Zach} erklärt nicht, was ihn davon abgehalten hat. Ein so einsichtsvoller Astronom hätte jedoch gewiss nicht auf einen Vorteil verzichtet, der die Wirkung vielleicht verdoppelt hätte, wenn das Gelände keine unüberwindlichen Hindernisse geboten hätte. Wäre die Wirkung des Mimet größer ausgefallen, hätten wir uns auch sehr eine vollständigere Kenntnis der räumlichen Abmessungen dieses Berges gewünscht. Unter den gegebenen Umständen wären die Schlüsse, die man daraus auf die relative Dichte des Minet und des ganzen Erdkörpers ziehen könnte, jedoch zu ungenau. Immerhin haben die Messungen von Herrn \textsc{von Zach}, indem sie wenigstens nachweisen, wie außergewöhnlich gering der Einfluss einer so bedeutenden Bergmasse ist, einen wichtigen Nutzen. Sie machen uns etwas misstrauischer gegenüber einer vorschnellen Berufung auf den möglichen Einfluss örtlicher Anziehungen, wenn die Messungen nicht zusammenpassen wollen.

Außer den genannten Untersuchungen enthält der sechste Abschnitt noch die Prüfung einiger Bestimmungen von \textsc{Cassini de Thury} in der Méridienne vérifiée, die von denen des Verfassers merklich abweichen. Viel größere, tatsächlich bis zur Entstellung gehende Fehler finden sich in \textsc{Cassini}s description géometrique de la France und in seiner großen Karte von Frankreich. --- Schließlich gibt der Verfasser noch einen Überblick über die Ergebnisse für die Positionen verschiedener Sterne, die teils aus seiner gegenwärtigen Arbeit, teils aus seinen eigenen früheren Beobachtungen und teils aus einer kritischen Erörterung der Beobachtungen anderer Astronomen folgen. Wir erwähnen hier nur die in vieler Hinsicht so wichtige Rektaszension des Polarsterns, für die \textsc{Mathieu}s Beobachtungen im Jahr 1812 eine Zunahme von 5 Zeitsekunden ergeben haben. Um \textsc{von Zach}s Bestimmung von 1790 damit in Übereinstimmung zu bringen, muss man eine Eigenbewegung voraussetzen. Herr \textsc{von Zach} findet sie durch den Vergleich mit \textsc{Lacaille}s Beobachtungen von 1750 auch vollständig bestätigt und setzt sie auf $+3''177$ im Bogen pro Jahr an. (Die genaueste Bestimmung der gegenwärtigen Rektaszension des Polarsterns werden uns die zahlreichen Beobachtungen liefern, die Herr \textsc{von Lindenau} seit fünf Jahren an der Seeberger Sternwarte angestellt hat. Von ihnen können wir schon vorläufig sagen, dass sie die von \textsc{Mathieu} gefundene Zunahme bestätigen.)

\Needspace{13\baselineskip}
Der begrenzte Raum unserer Zeitschrift erlaubt es uns nur, den Inhalt der übrigen Abschnitte von \textsc{von Zach}s Werk kurz zu streifen. Der siebte Abschnitt bestimmt die Höhe der Dreieckspunkte und einiger anderer Orte über dem Mittelmeer. Der achte Abschnitt gibt die Lage vieler Punkte in Marseille und der umliegenden Gegend nach ihrem Abstand vom Meridian und von der Senkrechten dazu an der dortigen Königlichen Sternwarte sowie nach ihrer Länge \SourceLeaf{579}und Breite an. Er enthält darüber hinaus mehrere weitere interessante kritische Untersuchungen über verschiedene Punkte in Marseille, die in der Geschichte der Astronomie bemerkenswert geworden sind. Schließlich gibt Herr \textsc{von Zach} in einem Anhang noch eine neue Berechnung sämtlicher astronomischer Beobachtungen, die \textsc{Maskelyne} am Shehallien angestellt hat. Das von \textsc{Maskelyne} selbst gefundene Endergebnis bleibt dabei unverändert.
